\section{Vertical Taxonomy}

Alignment.V1 uses one universal evidence and scoring framework across multiple
verticals. Vertical profiles describe differences in mandate, legal duties,
peer groups, data availability, and valid metrics; they do not create separate
scoring engines.

\subsection{Top-Level Families}

\begin{enumerate}
  \item \textbf{Legislative}: legislatures, councils, committees, and related
        public oversight bodies.
  \item \textbf{Executive}: departments, agencies, regulators, public
        administrators, health authorities, human services, and corrections.
  \item \textbf{Judicial}: judges, courts, clerks, prosecutors, and other
        adjudicative or justice-system entities.
  \item \textbf{Consumer-facing and commercial}: credit reporting agencies,
        lenders, insurers, employers, platforms, utilities, and other entities
        whose conduct directly affects consumers.
\end{enumerate}

These families are an organizational taxonomy, not a presumption that every
member has identical metrics. A Department of Corrections and a health
authority are both Executive entities but require different potential
profiles.

\subsection{First Validation Verticals}

The first cross-vertical validation set is deliberately heterogeneous:

\begin{center}
\begin{tabular}{p{4cm}p{5cm}p{5cm}}
\toprule
Family & Example profile & Primary validation purpose \\
\midrule
Executive & Department of Corrections & custody, safety, due process, rehabilitation, access, and accountability \\
Consumer-facing & Nationwide Consumer Reporting Agency & data accuracy, dispute handling, transparency, equity, and remediation \\
Judicial & Oregon circuit court or judge & procedural consistency, timeliness, access, and reasoned decision-making \\
Executive & Oregon Health Authority & public health administration, access, outcomes, and data stewardship \\
\bottomrule
\end{tabular}
\end{center}

Equifax is the first planned live entity instance, but the NCRA profile is the
reusable vertical artifact. A corrections department is the second validation
profile, not an Equifax-specific extension.
